\documentclass[aspectratio=169,10pt]{beamer}
\usetheme{default}
\setbeamertemplate{navigation symbols}{}
\usepackage{graphicx}\usepackage[T1]{fontenc}\usepackage{amsmath,amssymb}
\usepackage{listings}
\definecolor{ncsured}{HTML}{CC0000}\definecolor{ncsudk}{HTML}{990000}
\definecolor{ncsugray}{HTML}{6D6D6D}\definecolor{ncsuteal}{HTML}{008473}
\definecolor{codebg}{HTML}{F5F5F5}
\setbeamercolor{structure}{fg=ncsured}
\setbeamercolor{frametitle}{fg=white,bg=ncsured}
\setbeamercolor{title}{fg=ncsudk}
\setbeamercolor{itemize item}{fg=ncsured}
\setbeamercolor{itemize subitem}{fg=ncsured}
\setbeamerfont{frametitle}{series=\bfseries}
\setbeamertemplate{frametitle}{\nointerlineskip\vskip2pt%
  \begin{beamercolorbox}[wd=\paperwidth,ht=2.4ex,dp=1ex,leftskip=.3cm]{frametitle}%
  \small\bfseries\insertframetitle\end{beamercolorbox}}
\setbeamertemplate{footline}{\hbox{\begin{beamercolorbox}[wd=\paperwidth,ht=2.2ex,dp=1ex,leftskip=.3cm,rightskip=.3cm]{}%
  \tiny Optimal Market Making for Cryptocurrency Options \hfill Week 5 \hfill \insertframenumber\end{beamercolorbox}}}
\lstset{basicstyle=\ttfamily\scriptsize,backgroundcolor=\color{codebg},
  keywordstyle=\color{ncsured},commentstyle=\color{ncsugray},
  stringstyle=\color{ncsuteal},breaklines=true,columns=fullflexible,
  showstringspaces=false,frame=single,rulecolor=\color{ncsugray}}

\title{Week 5: The Avellaneda-Stoikov Engine}
\subtitle{Implementation, calibration, and backtest -- Milestone 1}
\date{North Carolina State University \,\textbar\, Summer 2026}
\titlegraphic{
\includegraphics[width=4.2cm]{/Users/dendisuhubdy/Github/ncsu/assets/ncsu_logo.png}}
\begin{document}
\begin{frame}[plain]\titlepage\begin{center}\tiny Industry Mentor: Dendi Suhubdy (Bitwyre)  |  Guest: Fenni Kang (Coincall)\end{center}\end{frame}
\begin{frame}{Contents}
\small\begin{enumerate}\setlength{\itemsep}{3pt}
  \item Learning objectives
  \item Lecture 1: From Closed Form to Code
  \item Lecture 2: Calibration \& Validation
  \item Illustration
  \item Derivation \& rationale
  \item Getting the data from Coincall
  \item Code: week5\_calibrate.py
  \item Glossary of key terms
  \item Reading
  \item References
\end{enumerate}
\end{frame}

\begin{frame}{Learning Objectives}
\begin{itemize}
  \item Implement the quotes as a tested PyTorch engine
  \item Calibrate the fill intensities by maximum likelihood
  \item Backtest on one month of Coincall BTC-PERP data
\end{itemize}
\end{frame}

\begin{frame}{Lecture 1: From Closed Form to Code}
\begin{itemize}
  \item Engine architecture; state from the L2 book
  \item Map continuous-time quantities to discrete events
  \item Calibrate lambda(delta) = A exp(-kappa delta) from fills
\end{itemize}
\end{frame}

\begin{frame}{Lecture 2: Calibration \& Validation}
\begin{itemize}
  \item Maximum-likelihood fit of A and kappa with goodness-of-fit
  \item Replay engine + fill simulator + P\&L attribution
  \item Reproducibility under a fixed seed
\end{itemize}
\end{frame}

\begin{frame}{Illustration}
\begin{center}\includegraphics[width=0.86\textwidth]{/Users/dendisuhubdy/Github/ncsu/weekly_presentations/figures/week05.png}\end{center}
\end{frame}

\section*{Derivation}
\begin{frame}{Derivation \& Rationale}
\[\mathcal{L}(A,\kappa)=\prod_i \lambda(\delta_i)\,e^{-\lambda(\delta_i)\,\Delta t_i}\]

\small\textbf{\color{ncsured}Why:} Fills are a (doubly stochastic) Poisson process, so the likelihood of the observed fills given quote distances $\delta_i$ is a product of Poisson densities. We fit by MAXIMUM LIKELIHOOD rather than least-squares because the data are event counts, not Gaussian residuals.
\end{frame}

\begin{frame}{Derivation \& Rationale}
\[-\ln\mathcal{L} = \sum_i\big[\lambda(\delta_i)\Delta t_i - \ln\lambda(\delta_i)\big]\]

\small\textbf{\color{ncsured}Why:} Minimizing this negative log-likelihood is convex in $\ln A$ and well-behaved in $\kappa$; PyTorch autograd gives the gradients, so a few hundred Adam steps converge.
\end{frame}

\begin{frame}{Getting the Data from Coincall}
\begin{itemize}
  \item Build the intensity-calibration set from trades vs. your quotes:
  \begin{itemize}\item \texttt{\small Trades:  GET /open/option/trade/lasttrade/v1/\{symbol\}  (or trade history)}\end{itemize}
  \begin{itemize}\item \texttt{\small Book:    GET /open/futures/order/orderbook/v1/\{symbol\}  for the mid}\end{itemize}
  \item Offline: coincall\_btc\_perp\_trades\_2025Q4.parquet + ...\_l2\_2025Q4.parquet
  \item A pre-aggregated coincall\_intensity\_calibration\_set.parquet is provided
  \item delta\_i = |fill price - mid at fill|; group fills by delta bucket for the MLE
\end{itemize}
\vspace{2pt}\tiny\textit{Cross-check exact paths at https://docs.coincall.com/}
\end{frame}

\begin{frame}[fragile]{Code --- week5\_calibrate.py}
\begin{lstlisting}
import torch
# delta: quote distances at which fills occurred; dt: exposure time per fill
delta = torch.tensor([...]); dt = torch.tensor([...])
logA  = torch.zeros(1, requires_grad=True)
kappa = torch.ones(1,  requires_grad=True)
opt = torch.optim.Adam([logA, kappa], lr=0.05)
for step in range(2000):
    lam = torch.exp(logA) * torch.exp(-kappa.clamp(min=1e-3) * delta)
    nll = (lam * dt - torch.log(lam)).sum()          # Poisson negative log-lik
    opt.zero_grad(); nll.backward(); opt.step()
print("A =", torch.exp(logA).item(), " kappa =", kappa.item())
\end{lstlisting}
\end{frame}
\begin{frame}{Glossary --- Key Terms}
\footnotesize\begin{description}
\item[\color{ncsured}Maximum likelihood (MLE)] fitting parameters by maximizing the likelihood of the data
\item[\color{ncsured}Poisson process] counting process for random arrivals (here, fills)
\item[\color{ncsured}Goodness-of-fit (chi-square)] a statistic measuring how well the fit matches the data
\item[\color{ncsured}Backtest] simulation of a strategy on historical data
\item[\color{ncsured}Fill simulator] model of which of the maker's quotes get executed
\item[\color{ncsured}P\&L attribution] decomposition of profit into its sources
\item[\color{ncsured}Reproducibility] identical results on re-run under a fixed random seed
\end{description}
\end{frame}

\begin{frame}{Reading}
\begin{itemize}
  \item Avellaneda-Stoikov (2008) [implementation]; Gueant (2016), Ch. 4
\end{itemize}
\end{frame}

\begin{frame}{References}
\scriptsize\begin{itemize}\setlength{\itemsep}{3pt}
  \item Avellaneda, M., \& Stoikov, S. (2008). High-frequency trading in a limit order book. Quantitative Finance, 8(3), 217-224.
  \item Gueant, O. (2016). The Financial Mathematics of Market Liquidity, Ch. 4. Chapman \& Hall/CRC.
  \item Cartea, Jaikumar, \& Penalva (2015). Algorithmic and High-Frequency Trading, Ch. 10.
  \item Coincall API documentation. https://docs.coincall.com/
\end{itemize}
\end{frame}

\end{document}